\chapter{Series Decomposition II}
\label{cap:filtros2}
% ── index ───────────────────────────────────────────────────────
\index{STL|textbf}
\index{stl@\texttt{stl()}|textbf}
\index{Christiano-Fitzgerald, filter of|textbf}
\index{christiano_fitzgerald@\texttt{christiano\_fitzgerald()}|textbf}
\index{loess}
\index{band-pass filter}
\index{seasonal decomposition}

% ─────────────────────────────────────────────────────────────────────────────
\section{STL and Christiano-Fitzgerald}

Ch.~\ref{cap:filtros} presented HP, BK and Holt-Winters. This chapter
covers two additional methods:

\begin{description}
  \item[STL] \emph{Seasonal-Trend decomposition via Loess} (Cleveland et al.~1990):
    decomposes $y_t = T_t + S_t + R_t$ iteratively via local regression (loess).
    Robust to outliers; allows time-varying seasonality.

  \item[Christiano-Fitzgerald (CF)] Asymmetric band-pass filter (2003):
    extracts frequencies in $[p_L, p_H]$ using weights from a
    non-symmetric window that adapts to the position in the series. Unlike the BK, it
    does not lose edge points — the cost is that weights change at each observation.
\end{description}

\subsection*{STL: iterative decomposition}

The STL algorithm alternates between:
\begin{enumerate}
  \item \emph{Seasonal loop}: smooths each seasonal position with loess
  \item \emph{Trend loop}: smooths the deseasonalized series with loess
\end{enumerate}
Convergence produces $T_t$ (trend), $S_t$ (seasonality) and $R_t$ (residual)
with robustness properties to outliers when \texttt{robust=true}.

\subsection*{Christiano-Fitzgerald Filter}

The CF applies weights $a_j$ satisfying $\sum_j a_j e^{i\omega j} \approx 1$
for $\omega \in [\omega_L, \omega_H]$ and $\approx 0$ outside the band.
For $t$ at the edges, the weights adjust to avoid using future observations.

% ─────────────────────────────────────────────────────────────────────────────
\section{DGP Verification}

DGP: monthly series with linear trend, annual cycle and semi-annual seasonality.
$T = 120$ months.

\begin{lstlisting}[caption={STL and Christiano-Fitzgerald}]
seed(42)
generate df trend  = 0.05 * t
generate df cycle  = 0.5 * sin(t * 3.14159 / 12)
generate df season = 0.3 * sin(t * 3.14159 / 6)
generate df y      = 10.0 + trend + cycle + season + noise

stl(df, y, period=12)
display df

christiano_fitzgerald(df, y, low=6, high=32)
display df
\end{lstlisting}

\subsection*{STL}

\begin{lstlisting}[language={}, basicstyle=\ttfamily\small, frame=none,
  numbers=none, backgroundcolor=\color{codbg}]
Seasonal Decomposition (STL)
  Observations: 120
  Trend mean:   13.0243
  Seasonal std:  0.1943
  Residual std:  0.1868
\end{lstlisting}

The STL adds columns \texttt{y\_trend}, \texttt{y\_seasonal} and \texttt{y\_resid}
to the DataFrame. Residual std $< 0{,}2$ — the true noise ($\sigma = 0{,}3$) is
well separated from the other components.

\subsection*{Christiano-Fitzgerald}

\begin{lstlisting}[language={}, basicstyle=\ttfamily\small, frame=none,
  numbers=none, backgroundcolor=\color{codbg}]
cffilter: periods [6,32]  →  y_cycle added to df
  t=1:  y_cycle = 1.04
  t=6:  y_cycle = 1.19    (dominant annual cycle)
  t=12: y_cycle = 0.87
\end{lstlisting}

The CF retains periodic variations between 6 and 32 months (1{,}5 to 8 years — business
cycles). The advantage over the BK is that all 120 observations are kept
(no edge trimming).

% ─────────────────────────────────────────────────────────────────────────────
\section{Syntax}

\begin{lstlisting}
stl(df, var, period=12)                        // adds trend, seasonal, resid
christiano_fitzgerald(df, var, low=6, high=32) // adds var_cycle
\end{lstlisting}

\begin{center}
\begin{tabular}{llll}
\toprule
\textbf{Filter} & \textbf{Type} & \textbf{Edge loss} & \textbf{Seasonality} \\
\midrule
HP              & low-pass      & Yes (endpoints)       & No \\
BK              & band-pass     & Yes ($k$ each side)   & No \\
CF              & band-pass     & No                    & No \\
STL             & decomposition & No                    & Yes \\
\bottomrule
\end{tabular}
\end{center}

\begin{nota}
  The function \texttt{hamilton} in Hayashi refers to the Markov
  Switching model of Hamilton (1989), \emph{not} to the Hamilton filter (2018)
  (ch.~\ref{cap:regime}). The Hamilton 2018 filter — which proposes replacing
  the HP with an OLS regression with 4 lags and horizon $h = 8$ —
  is not implemented separately.
\end{nota}
